\section{有理线性泛函}

\begin{itemize}
	\item $K$是除环,$K^{\prime}$是其子除环,$\rho:K^{\prime}\hookrightarrow K$是典范嵌入.
	\item $V$是右$K$-向量空间,其$K^{\prime}$-结构是$V^{\prime}$
\end{itemize}

\begin{definition}{有理线性泛函}{}
	设$f\in V^{\vee}$.若$\rho_{\ast}\left(f\right)\left(V^{\prime}\right)\subseteq K_{d}^{\prime}$,则称$f$是$K^{\prime}$-有理线性泛函.$V$上的$K^{\prime}$-有理线性泛函组成的集合记作$\Rat_{K^{\prime}}\left(V,V^{\prime}\right)$.
\end{definition}

\begin{proposition}{}{}
	设$\Gamma:\left(V^{\prime}\right)^{\vee}\to\rho_{\ast}\left(V^{\vee}\right)$是典范映射.
	\begin{enumerate}
		\item $\Gamma$是单射,并且$\im\left(\Gamma\right)=\Rat_{K^{\prime}}\left(V,V^{\prime}\right)$.
		\item $\Rat\left(V,V^{\prime}\right)$的自由子集都是$V^{\vee}$的自由子集.
	\end{enumerate}
\end{proposition}

\begin{proposition}{}{}
	若$V$是有限维的,则$\Rat_{K^{\prime}}\left(V,V^{\prime}\right)$是$V$的$K^{\prime}$-结构.
\end{proposition}

\begin{proposition}{}{}
	设$W$是$V$的子空间,则$W$是$V$的$K^{\prime}$-有理子空间$\iff$存在$H\subseteq\Rat_{K^{\prime}}\left(V,V^{\prime}\right)$,使得$W=H^{\circ}$.
\end{proposition}
